% ============================================================
%  第二章  推断因果性的理论
% ============================================================
\chapter{推断因果性的理论}

\begin{flushright}
\kaishu\small
我宁愿发现一条因果律，也不愿做波斯之王。\\[2pt]
——德谟克利特（Democritus，公元前 460—370）
\end{flushright}

\section*{序言}
\addcontentsline{toc}{section}{序言}

从原始数据中学习因果关系的可能性，自休谟（1711—1776）时代起就列在哲学家的梦想清单上。20 世纪 80 年代中期，随着图与概率依赖之间的数学关系被揭示，这种可能性进入了形式化处理与可行计算的领域。本章所述进路是 Rebane and Pearl (1987) 与 Pearl (1988b, 第 8 章) 的产物；后者说明，只要对数据生成过程作出某些假设（例如假定它具有树结构），就可以从非时间的统计数据中推断因果关系。从更弱的结构假设（例如一般的\emph{有向无环图}）出发推断因果关系的前景，促成了三所大学——加州大学洛杉矶分校（UCLA）、卡内基梅隆大学（CMU）与斯坦福大学——的平行研究。UCLA 与 CMU 的团队采取的进路是：在数据中搜寻条件独立模式，这些模式揭示了底层结构的片段，再把这些片段拼装成自洽的因果模型（或此类模型的集合）。斯坦福团队采取贝叶斯进路，用数据更新赋予候选因果结构的先验概率（Cooper and Herskovits 1991）。UCLA 与 CMU 的工作导致了相似的理论和几乎相同的发现算法，并在 TETRAD II 程序中实现（Spirtes et al. 1993）。贝叶斯进路此后为若干研究团队所发展（Singh and Valtorta 1995; Heckerman et al. 1994），如今已成为若干基于图的学习方法的基础（Jordan 1998）。本章叙述 Tom Verma 与我在 1988—1992 年间所走的进路，并简要总结 CMU 团队及其他学者提出的相关推广、精炼与改进。这一发展背后的某些哲学理据，主要是极小性假设，同样隐含在贝叶斯进路之中（\S 2.9.1）。

把原因发现自动化的基本想法——以及这一想法在计算机程序中的具体实现——曾在多个论坛上引起激烈争论（Cartwright 1995a; Humphreys and Freedman 1996; Cartwright 1999; Korb and Wallace 1997; McKim and Turner 1997; Robins and Wasserman 1999）。这场争论的某些方面将在本章末的讨论节（\S 2.9.1）中处理。

本章承认统计关联在逻辑上并不蕴含因果关系，转而追问二者之间是否存在较弱的联系。具体而言，我们要问：
\begin{enumerate}[label=\arabic*.]
\item 是什么线索促使人们在不受控的观测中感知到因果关系？
\item 需要什么假设，才能使我们从这些线索推断出因果模型？
\item 所推断出的模型，能否就支撑观测的因果机制告诉我们某些有用的东西？
\end{enumerate}

\S 2.2 定义因果模型与因果结构的概念，并把因果发现的任务描述为科学家与自然之间的一场归纳博弈。\S 2.3 通过引入“极小模型”语义学——奥卡姆剃刀的语义版本——把这场归纳博弈形式化，并举例说明：与通常的流行说法相反，遵循这一归纳推理的标准规范，因果关系是可以与虚假共变区分开来的。\S 2.4 指出一个称为\emph{稳定性}（或忠实性）的条件，在该条件下存在有效算法，能够揭示此处所定义的因果影响结构。\S 2.5 引入的 IC 算法假定所有变量均被观测，可揭示与数据相容的一切因果模型之集合。\S 2.6 描述的 IC* 算法则表明，当某些变量不可观测时，它仍能揭示许多（虽非全部）有效的因果关系。\S 2.7 从 IC* 算法中提炼出推断因果影响所依据的本质条件，并把它们作为真正影响与虚假关联的独立定义提出，分别处理有、无时间信息两种情形。\S 2.8 对因果性在时间方面与统计方面何以令人困惑地普遍一致，给出一种解释。最后，\S 2.9 总结本章的主张，重新阐明导致这些主张的假设，并结合正在进行的争论为这些假设提供新的辩护。

\section{引论——基本直觉}

一个自主的智能系统若要为自身环境建立可用的模型，就不能完全依赖预编程的因果知识；它必须能够把直接观测转译为因果关系。然而，既然统计分析是由共变而非因果驱动的，又假定人类知识的大部分来自被动观测，我们就仍须找出促使人们在数据中感知因果关系的线索，并且必须找到一种能够模拟这种感知的计算模型。

时间在先通常被认为是定义因果关系的必要条件，而且它无疑是人们用以区分因果关联与其他类型关联的最重要线索之一。据此，大多数因果理论都明确要求原因在时间上先于结果（Reichenbach 1956; Good 1961; Suppes 1970; Shoham 1988）。然而，仅凭时间信息无法把真正的因果关系与由未知因素造成的虚假关联区分开来——气压计在雨前下降，却不导致下雨。事实上，统计学与哲学文献一直严厉警告分析者：除非事先知道一切因果相关因素，或者能够审慎地操纵某些变量，否则不可能作出真正的因果推断（Fisher 1935; Skyrms 1980; Cliff 1983; Eells and Sober 1983; Holland 1986; Gardenfors 1988; Cartwright 1989）。\footnote{一些流行的说法包括：“没有操纵就没有因果”（Holland 1986）；“输入无因，输出无因”（Cartwright 1989）；“没有哪个计算机程序能考虑不在分析之中的变量”（Cliff 1983）。我最喜欢的一句是：“因果第一，操纵第二。”}这两种条件在正常的学习环境中都无法实现，于是问题依然是：因果知识究竟是如何从经验中获得的？

本章所探讨的线索，来自某些具有因果组织特征的统计关联模式——这些模式事实上只有用因果方向性才能得到有意义的解释。例如，考虑三个事件之间如下非传递的依赖模式：$A$ 与 $B$ 相依，$B$ 与 $C$ 相依，而 $A$ 与 $C$ 独立。若请一个人为这样三个事件举例，他所举的例子必定把 $A$ 与 $C$ 描绘为两个独立的原因、把 $B$ 描绘为它们的共同结果，即
\[
A\to B\leftarrow C .
\]
（我最喜欢的例子是：$A$ 与 $C$ 是两枚公平硬币的结果，$B$ 表示一只铃，只要任一硬币正面朝上它就响。）用“$B$ 是原因、$A$ 与 $C$ 是结果”的情景去拟合这一依赖模式，在数学上可行，却非常不自然（鼓励读者一试）。

这类思想实验告诉我们：某些依赖模式在缺乏时间信息的情况下，在概念上就带有某些因果方向性的特征，而不带另一些方向性的特征。Reichenbach (1956) 最先追问这些模式的起源，他认为它们是自然的特征，反映了热力学第二定律。Rebane and Pearl (1987) 反过来提出问题：三种基本因果子结构
\[
X\to Y\to Z,\qquad X\leftarrow Y\to Z,\qquad X\to Y\leftarrow Z
\]
所对应的依赖之间的区别，能否用来揭示底层数据生成过程中真正的因果影响？他们很快意识到，判定 $X$ 与 $Y$ 之间因果关系方向的关键，在于“存在第三个变量 $Z$，它与 $Y$ 相关而不与 $X$ 相关”，正如碰撞结构 $X\to Y\leftarrow Z$ 那样；据此他们发展出一个算法，能在他们所考虑的因果图类（即多树）中同时恢复边与方向。

本章的探究把这些直觉形式化，并把 Rebane–Pearl 恢复算法推广到一般图，包括含有未观测变量的图。

\section{因果发现框架}

我们把因果发现的任务视为科学家与自然之间的一场归纳博弈。自然拥有稳定的因果机制；在细致的描述层面上，这些机制是变量之间的确定性函数关系，其中有些变量不可观测。这些机制以无环结构的形式组织起来，科学家则试图从可得的观测中识别出这一结构。

\begin{definition}[因果结构]
变量集 $V$ 的\emph{因果结构}是一个有向无环图（DAG），其中每个节点对应 $V$ 中的一个不同元素，每条连线表示相应变量之间的一个直接函数关系。
\end{definition}

因果结构是构造“因果模型”的蓝图——因果模型精确地规定了每个变量如何受其在 DAG 中父节点的影响，如 \textup{(1.40)} 的结构方程模型那样。这里我们假定，自然可以自由地在每个结果与其原因之间施加任意的函数关系，然后再通过引入任意（但相互独立）的扰动来扰动这些关系。这些扰动反映了自然以某种未公开的概率函数所支配的“隐藏”或不可测的条件。

\begin{definition}[因果模型]
\emph{因果模型}是一个二元组 $M=\langle D,\Theta_D\rangle$，由因果结构 $D$ 与一组与 $D$ 相容的参数 $\Theta_D$ 构成。参数 $\Theta_D$ 为每个 $X_i\in V$ 指定一个函数 $x_i=f_i(\pa_i,u_i)$，并为每个 $u_i$ 指定一个概率测度 $P(u_i)$；其中 $\PA_i$ 是 $X_i$ 在 $D$ 中的父节点，每个 $U_i$ 是按 $P(u_i)$ 分布、且与其余一切 $u$ 独立的随机扰动。
\end{definition}

正如第一章（定理 1.4.1）所见，扰动相互独立的假设使模型成为马尔可夫的，即每个变量在给定其在 $D$ 中的父节点时，与它的所有非后代独立。马尔可夫假设在人类话语中的普遍性，或许反映了我们认为有助于理解自然的那些模型的粒度。我们可以从确定性的极端出发：那里一切变量都被微观地详尽阐明，马尔可夫条件当然成立。当我们通过聚合变量、引入概率以概括被省略的变量而上升到宏观抽象时，就需要决定抽象到哪一步便已过头、有用的因果性质在何处丧失。显然，我们的祖先（我们因果思想的作者们）已把马尔可夫条件认作在这种抽象中值得保护的性质：不由共同原因解释的相关被视为虚假的，含有此类相关的模型被视为不完整的。马尔可夫条件指引我们判断一个父节点集 $\PA_i$ 何时算作完整——即它包含了变量 $X_i$ 的一切相关直接原因。它允许我们把其中某些原因留在 $\PA_i$ 之外（用概率加以概括），但前提是这些原因不影响系统中被建模的其他变量。若模型中的 $\PA_i$ 过窄，就会出现同时影响若干变量的扰动项，马尔可夫性便丧失了。这类扰动将被显式地当作“潜”变量处理（见定义 2.3.2）。一旦我们承认潜变量的存在、并把它们的存在显式表示为图中的节点，马尔可夫性便得以恢复。

因果模型 $M$ 一经形成，就在系统变量上定义了一个联合概率分布 $P(M)$。该分布反映了因果结构的某些特征（例如，每个变量在给定其父节点取值时必与其祖父节点独立）。随后，自然允许科学家考察一个选定的“被观测”变量子集 $O\subseteq V$，并就 $P[O]$（$O$ 上的概率分布）提出问题，却把底层的因果模型连同因果结构一并隐藏起来。我们所要探究的，是从概率分布 $P[O]$ 的特征中恢复 DAG 拓扑 $D$ 的可行性。\footnote{这一表述对科学发现的实际任务作了若干理想化。例如，它假定科学家直接获得分布，而不是从分布中抽样得到的事件。此外，我们假定被观测变量确实出现在原初因果模型中，而不是其某种聚合。聚合可能导致反馈环，本章不予讨论。}

\section{模型偏好（奥卡姆剃刀）}

原则上，由于 $V$ 未知，能够拟合给定分布的模型数目是无界的：每个模型都诉诸一组不同的“隐藏”变量，并以不同的因果关系连接被观测变量。因此，若对所考虑的模型类型不加限制，科学家便无法就现象背后的结构作出任何有意义的断言。例如，每个概率分布 $P[O]$ 都可以由一个结构生成：其中没有任何被观测变量是另一个的原因，而是所有变量都是某个潜在共同原因 $U$ 的后果。\footnote{这可以通过如下方式实现：令 $U$ 的状态数与 $O$ 相同，给 $U$ 赋予先验分布 $P(u)=P(o(u))$（其中 $o(u)$ 是 $O$ 中对应于状态 $u$ 的单元），并令每个被观测变量 $O_i$ 取 $o(u)$ 中相应的值。}同样，假定 $V=O$ 却缺乏时间信息时，科学家永远无法排除这样一种可能：底层结构是一个完全的、无环的、次序任意的图——这种结构（在参数选择恰当时）可以模拟任何模型的行为，而不论变量次序如何。然而，遵循科学归纳的标准规范，凡当我们能找到更简单、更不繁复而又同样与数据相容的理论时，就有理由排除原有的理论（见定义 2.3.5）。经受住这一筛选过程的理论称为\emph{极小}的。借助这一概念，我们可以如下构造“推断因果性”的（初步）定义。

\begin{definition}[推断因果性（初步）]
若在一切与数据相容的极小结构中，都存在一条从 $X$ 到 $Y$ 的有向路径，则称变量 $X$ 对变量 $Y$ 有\emph{因果影响}。
\end{definition}

这里我们把因果结构等同于科学理论，因为二者都包含一组可调整以拟合数据的自由参数。我们之所以把定义 2.3.1 视为初步的，是因为它假定所有变量都被观测。接下来的几个定义把极小性概念推广到含有未观测变量的结构。

\begin{definition}[潜结构]
\emph{潜结构}是一个二元组 $L=\langle D,O\rangle$，其中 $D$ 是 $V$ 上的因果结构，$O\subseteq V$ 是一个被观测变量集。
\end{definition}

\begin{definition}[结构偏好]
一个潜结构 $L=\langle D,O\rangle$ 优于另一个潜结构 $L'=\langle D',O\rangle$（记作 $L\preceq L'$），当且仅当 $D'$ 能在 $O$ 上模拟 $D$——即当且仅当对每个 $\Theta_D$ 都存在 $\Theta'_{D'}$，使得
\[
P[O](\langle D',\Theta'_{D'}\rangle)=P[O](\langle D,\Theta_D\rangle).
\]
两个潜结构\emph{等价}，记作 $L\equiv L'$，当且仅当 $L\preceq L'$ 且 $L'\preceq L$。\footnote{我们用简洁的说法“优于”表示“优于或等价于”，这一关系也被称为“是……的子模型”。}
\end{definition}

注意，定义 2.3.3 所施加的简单性偏好，是以结构的表达力、而不是以其句法描述来衡量的。例如，一个潜结构 $L_1$ 所诉诸的参数可能远多于 $L_2$，但只要 $L_2$ 能容纳更丰富的可观测分布集，$L_1$ 仍会被偏好。科学家偏好较简单的理论，一个原因是这类理论约束更强、因而更可否证；它们使科学家较少有机会“事后”过度拟合数据，因此一旦拟合成功便具有更大的可信度（Popper 1959; Pearl 1978; Blumer et al. 1987）。

我们还注意到，因果结构所蕴含的独立关系集对其表达力——即模拟其他结构的能力——施加了限制。事实上，只要存在一个 $L_1$ 允许而 $L_2$ 禁止的可观测依赖，$L_1$ 就不能优于 $L_2$。因此，偏好与等价的检验有时可以化归为对诱导依赖的检验，而后者又可直接由 DAG 的拓扑确定，完全不必理会参数集。在没有隐藏变量时正是如此（见定理 1.2.8），但在一切潜结构中并不普遍成立。Verma and Pearl (1990) 表明，某些潜结构对被观测分布施加的是数值约束而非独立约束（例如见 \S 8.4 式 \textup{(8.21)}—\textup{(8.23)}）；这使检验模型偏好的任务变得复杂，但仍允许我们把推断因果性的语义定义（定义 2.3.1）推广到潜结构。

\begin{definition}[极小性]
潜结构 $L$ 相对于潜结构类 $\mathcal{L}$ 是\emph{极小}的，当且仅当 $\mathcal{L}$ 中没有任何成员严格优于 $L$——即当且仅当对每个 $L'\in\mathcal{L}$，只要 $L\preceq L'$ 就有 $L'\preceq L$。
\end{definition}

\begin{definition}[相容性]
潜结构 $L=\langle D,O\rangle$ 与 $O$ 上的分布 $\hat{P}$ \emph{相容}，当且仅当 $D$ 能容纳某个生成 $\hat{P}$ 的模型——即存在参数化 $\Theta_D$ 使得 $P[O](\langle D,\Theta_D\rangle)=\hat{P}$。
\end{definition}

显然，$L$ 与 $\hat{P}$ 相容的一个必要（有时也是充分）条件是：$L$ 能解释 $\hat{P}$ 中所体现的一切依赖。

\begin{definition}[推断因果性]
给定 $\hat{P}$，变量 $C$ 对变量 $E$ 有\emph{因果影响}，当且仅当在一切与 $\hat{P}$ 相容的极小潜结构中，都存在一条从 $C$ 到 $E$ 的有向路径。
\end{definition}

我们认为这个定义是规范性的，因为它立足于科学研究中最少争议的规范之一：以语义方式表述的奥卡姆剃刀。然而，与任何科学探究一样，我们并不声称这个定义必定总能识别出自然中稳定的物理机制。它所识别的是我们能够从非实验数据中合理推断出的机制；不仅如此，它还保证：任何替代机制都比所推断的机制更不可信，因为替代者将需要对参数（即函数）作更牵强、更事后的调整才能拟合数据。

作为定义 2.3.6 所识别出的一个因果关系的例子，设想在四个变量 $\{a,b,c,d\}$ 上所作的观测揭示出两个独立关系：“$a$ 与 $b$ 独立”以及“给定 $c$ 时 $d$ 与 $\{a,b\}$ 独立”。再假定数据并未揭示除这两个关系之逻辑后承之外的任何其他独立。这种依赖模式典型地对应于下列变量：$a=$ 患感冒，$b=$ 患花粉症，$c=$ 想打喷嚏，$d=$ 想擦鼻子。

不难看出，图 2.1 中的结构 (a) 与 (b) 是极小的，因为它们恰好蕴含所观测到的独立关系而无其他。\footnote{为验证 (a) 与 (b) 等价，注意若令连线 $a\leftarrow *$ 施加两变量相等这一约束，(b) 便可模拟 (a)；反之，(a) 也能模拟 (b)，因为 (a) 能生成具有 (b) 所蕴含之独立关系的每个分布。（关于从图中“读出”条件独立关系的理论与方法，见 \S 1.2.3 或 Pearl 1988b。）}此外，任何用从 $d$ 指向 $c$ 的箭头、或用二者之间的隐藏共同原因（$*$）来解释 $c$ 与 $d$ 之间观测依赖的结构，都不可能是极小的；因为任何这样的结构都能“胜过”图 2.1(a)（或图 2.1(b)）所示的结构，而后者反映了全部观测独立。例如，图 2.1(c) 的结构不同于图 2.1(a)，它能容纳 $a$ 与 $b$ 之间关系任意的分布。同样，图 2.1(d) 不是极小的，因为它未能施加给定 $c$ 时 $d$ 与 $\{a,b\}$ 之间的条件独立，从而会容纳那些给定 $c$ 时 $d$ 与 $\{a,b\}$ 相依的分布。相反，图 2.1(e) 与数据不相容，因为它施加了一个未被观测到的、$\{a,b\}$ 与 $d$ 之间的边缘独立。

\begin{figure}[htbp]
\centering
\begin{tikzpicture}[causal,scale=0.92,every node/.style={transform shape}]
  % (a)
  \begin{scope}
    \node[circle,fill,inner sep=1.3pt] (a) at (-0.9,0.9) {}; \node[font=\small,left=0.5mm] at (-0.9,0.9) {$a$};
    \node[circle,fill,inner sep=1.3pt] (b) at (0.9,0.9) {};  \node[font=\small,right=0.5mm] at (0.9,0.9) {$b$};
    \node[circle,fill,inner sep=1.3pt] (c) at (0,0) {};      \node[font=\small,right=0.6mm] at (0,0) {$c$};
    \node[circle,fill,inner sep=1.3pt] (d) at (0,-1.0) {};   \node[font=\small,below=0.4mm] at (0,-1.0) {$d$};
    \draw[ed] (a)--(c); \draw[ed] (b)--(c); \draw[ed] (c)--(d);
    \node[font=\small] at (0,-1.75) {(a)};
  \end{scope}
  % (b)
  \begin{scope}[xshift=3.4cm]
    \node[circle,fill,inner sep=1.3pt] (a) at (-0.9,0.55) {}; \node[font=\small,left=0.5mm] at (-0.9,0.55) {$a$};
    \node[circle,fill,inner sep=1.3pt] (b) at (0.9,0.9) {};   \node[font=\small,right=0.5mm] at (0.9,0.9) {$b$};
    \node[circle,fill,inner sep=1.3pt] (s) at (0,1.5) {};     \node[font=\small,above=0.3mm] at (0,1.5) {$*$};
    \node[circle,fill,inner sep=1.3pt] (c) at (0,0) {};       \node[font=\small,right=0.6mm] at (0,0) {$c$};
    \node[circle,fill,inner sep=1.3pt] (d) at (0,-1.0) {};    \node[font=\small,below=0.4mm] at (0,-1.0) {$d$};
    \draw[ed] (s)--(a); \draw[ed] (s)--(c); \draw[ed] (b)--(c); \draw[ed] (c)--(d);
    \node[font=\small] at (0,-1.75) {(b)};
  \end{scope}
  % (c)
  \begin{scope}[xshift=7.0cm]
    \node[circle,fill,inner sep=1.3pt] (a) at (-0.9,0.9) {}; \node[font=\small,left=0.5mm] at (-0.9,0.9) {$a$};
    \node[circle,fill,inner sep=1.3pt] (b) at (0.9,0.9) {};  \node[font=\small,right=0.5mm] at (0.9,0.9) {$b$};
    \node[circle,fill,inner sep=1.3pt] (c) at (0,0) {};      \node[font=\small,right=0.6mm] at (0,0) {$c$};
    \node[circle,fill,inner sep=1.3pt] (d) at (0,-1.0) {};   \node[font=\small,below=0.4mm] at (0,-1.0) {$d$};
    \draw[ed] (a)--(b); \draw[ed] (a)--(c); \draw[ed] (b)--(c); \draw[ed] (c)--(d);
    \node[font=\small] at (0,-1.75) {(c)};
  \end{scope}
  % (d)
  \begin{scope}[yshift=-3.5cm,xshift=1.0cm]
    \node[circle,fill,inner sep=1.3pt] (a) at (-0.7,0.9) {}; \node[font=\small,above=0.4mm] at (-0.7,0.9) {$a$};
    \node[circle,fill,inner sep=1.3pt] (b) at (0.7,0.9) {};  \node[font=\small,above=0.4mm] at (0.7,0.9) {$b$};
    \node[circle,fill,inner sep=1.3pt] (c) at (0,0) {};      \node[font=\small,right=0.6mm] at (0,0) {$c$};
    \node[circle,fill,inner sep=1.3pt] (d) at (0,-1.1) {};   \node[font=\small,below=0.4mm] at (0,-1.1) {$d$};
    \node[circle,fill,inner sep=1.3pt] (s1) at (-2.0,0.35) {}; \node[font=\small,left=0.6mm] at (-2.0,0.35) {$*$};
    \node[circle,fill,inner sep=1.3pt] (s2) at (2.0,0.35) {};  \node[font=\small,right=0.6mm] at (2.0,0.35) {$*$};
    \draw[ed] (s1)--(a); \draw[ed] (s1)--(c); \draw[ed] (s2)--(b); \draw[ed] (s2)--(c);
    \draw[ed] (a)--(c); \draw[ed] (b)--(c);
    \draw[ed] (s1)--(d); \draw[ed] (s2)--(d);
    \node[font=\small] at (0,-1.85) {(d)};
  \end{scope}
  % (e)
  \begin{scope}[yshift=-3.5cm,xshift=5.6cm]
    \node[circle,fill,inner sep=1.3pt] (a) at (-0.7,0.9) {}; \node[font=\small,above=0.4mm] at (-0.7,0.9) {$a$};
    \node[circle,fill,inner sep=1.3pt] (b) at (0.7,0.9) {};  \node[font=\small,above=0.4mm] at (0.7,0.9) {$b$};
    \node[circle,fill,inner sep=1.3pt] (c) at (0,0) {};      \node[font=\small,right=0.6mm] at (0,0) {$c$};
    \node[circle,fill,inner sep=1.3pt] (d) at (0,-1.1) {};   \node[font=\small,below=0.4mm] at (0,-1.1) {$d$};
    \node[circle,fill,inner sep=1.3pt] (s) at (1.8,0) {};    \node[font=\small,right=0.6mm] at (1.8,0) {$*$};
    \draw[ed] (a)--(c); \draw[ed] (b)--(c); \draw[ed] (s)--(c); \draw[ed] (s)--(d);
    \node[font=\small] at (0,-1.85) {(e)};
  \end{scope}
\end{tikzpicture}
\caption{说明 (a) 与 (b) 的极小性、以及推断关系 $c\to d$ 之理据的因果结构。节点（$*$）表示一个状态数任意的隐藏变量。}
\label{fig:2.1}
\end{figure}

这个例子（取自 Pearl and Verma 1991）揭示了因果性与概率之间一个引人注目的联系：某些概率依赖模式（本例中，除 $a\perp b$ 与 $(d\perp\{a,b\}\mid c)$ 之外的一切依赖）蕴含着明确的因果依赖（本例中为 $c\to d$），而对潜变量之有无不作任何假设。\footnote{因果性的标准概率定义（如 Suppes 1970; Eells 1991）总是要求知道一切可能影响被观测变量的相关因素（见 \S 7.5.3）。}这一蕴含所援引的唯一假设就是极小性——过度拟合数据的模型被排除在外。
